\documentclass[11pt]{article}
\usepackage[a4paper,margin=24mm]{geometry}
\usepackage[T1]{fontenc}
\usepackage{lmodern,microtype}
\usepackage{amsmath,amssymb,amsthm,mathtools}
\usepackage{booktabs,array,longtable,needspace}
\usepackage{xurl}
\usepackage[hidelinks]{hyperref}
\setlength{\parindent}{0pt}
\setlength{\parskip}{5pt}
\setlength{\emergencystretch}{2em}
\renewcommand{\arraystretch}{1.16}
\newtheorem{theorem}{Theorem}[section]
\newtheorem{lemma}[theorem]{Lemma}
\newtheorem{proposition}[theorem]{Proposition}
\newtheorem{corollary}[theorem]{Corollary}
\theoremstyle{remark}
\newtheorem{remark}[theorem]{Remark}
\newcommand{\op}{\mathrm{op}}
\newcommand{\norm}[1]{\left\lVert#1\right\rVert_{\op}}
\newcommand{\Cstar}{C_{\mathrm{normal}}}
\newcommand{\dd}{\,\mathrm{d}}
\newcommand{\ii}{\mathrm{i}}
\newcommand{\diag}{\operatorname{diag}}
\newcommand{\ran}{\operatorname{ran}}
\newcommand{\esssup}{\operatorname*{ess\,sup}}
\newcommand{\Rformal}{\mathcal R}
\begin{document}
\begin{center}
{\LARGE\bfseries SP-07: Exact completion barriers\\[4pt] and a two-line variational model}
\\[10pt] Sidney Holden\\[3pt]{\small Center for Computational Biology, Flatiron Institute, Simons Foundation}
\end{center}

\textbf{Status.} The unrestricted sharp constant is not determined. This continuation solves a reflection-completion relaxation for every odd regular polygon and solves the phase optimization in a stationary two-line model. It does not improve the inherited finite lower bound or the literature upper bound. In particular, the number near $1.0373054$ below is \emph{not} a new bound for SP-07.

\begin{abstract}
For a regular odd polygon of selected spectral values, we obtain an exact formula for the least anticommutator norm over all Hermitian-unitary completions of a prescribed circulant overlap. A norm-preserving orbit dilation removes the circulant assumption from the resulting obstruction: whenever the selected compression has an eigenvalue $+1$ or $-1$, the selected pair-sum separation cannot exceed the full anticommutator norm. Thus this geometry cannot produce a greater-than-one selected Hall quotient, in arbitrary completion dimension.

An explicit sequence of compressed optima has quotient tending to $\pi/2$, while its best completed selected quotient tends to zero. Separately, a bounded anticommutator of an unbounded normal two-line operator reduces to a scalar speed-control problem with an exact integral solution. Rational arithmetic certifies its formal selected quotient at one parameter pair between $1.0373054176426$ and $1.0373054176428$. An explicit infinite matching, however, has a ratio below one. This identifies the finite Hall deficit and its boundary realization as essential missing ingredients, rather than silently transferring an infinite or compressed optimum to SP-07.
\end{abstract}

\section{The target and the inherited baseline}
For normal matrices $A,B\in\mathbb C^{n\times n}$, list their eigenvalues with multiplicity as $\alpha_i,\beta_i$ and put
\[
 d_\infty(A,B)=\min_{\pi\in S_n}\max_i|\alpha_i-\beta_{\pi(i)}|,
 \qquad
 \Cstar=\sup_{n\ge2}\sup_{A\ne B\text{ normal}}
 \frac{d_\infty(A,B)}{\norm{A-B}}.
\]
SP-07 asks for one sharp constant valid in all dimensions, with a matching construction or limiting sharpness argument \cite{sp07}. Neither a sharp subclass result nor an optimum of a relaxation answers that question.

The preceding package proves the exact rank-one-reflection constant $\sqrt{28/27}$ and a finite-order reduction $n\le 3r^2+7r-1$ at each fixed negative reflection rank $r$ \cite{prior4}. It does not give a rank cutoff. The earlier real-reflection embedding shows that arbitrary normal pairs are represented, without changing their ratio, within real reflection anti-pairs after increasing dimension. These inherited results remain distinct from the new completion calculations below.

The strongest inherited rational example has order $193$, matching distance exactly $2$, and perturbation norm strictly below $19281027/10^7$. It has been rerun in this continuation by both of its exact positivity methods. With the upper bound reported in Tang's versioned preprint \cite{tang}, the bracket remains
\begin{equation}\label{eq:bracket}
 \boxed{\quad \frac{20000000}{19281027}<\Cstar<2.9038872828.\quad}
\end{equation}
Only the lower endpoint is backed here by a retained finite arithmetic certificate; the upper endpoint is a cited literature bound, not rederived. No new finite lower endpoint is claimed.

\section{An exact fixed-overlap completion formula}
Fix an odd integer $p\ge3$, a radius $R>0$, and $\omega=e^{2\pi\ii/p}$. The selected normal block is
\[
 D_s=R\diag(1,\omega,\ldots,\omega^{p-1}).
\]
Its selected pair-sum separation is
\begin{equation}\label{eq:polygon-gap}
 \delta_p=\min_{i,j}|(D_s)_{ii}+(D_s)_{jj}|
 =2R\sin\frac{\pi}{2p}.
\end{equation}
Indeed, the relevant moduli are $2R|\cos(\pi(i-j)/p)|$, and $p$ is odd.

Let $X$ be a Hermitian circulant contraction on the selected space, with $X\mathbf1=\mathbf1$. For the unitary Fourier matrix
\[
 F_{jk}=p^{-1/2}\omega^{jk},\qquad 0\le j,k<p,
\]
write $F^*XF=\diag(\lambda_0,\ldots,\lambda_{p-1})$, where $\lambda_0=1$ and $-1\le\lambda_j\le1$. In this basis $F^*D_sF=RT$, where $T a_j=a_{j+1\bmod p}$.

We allow a Hermitian-unitary completion $U$ of $X$ in any finite dimension, and a block-diagonal extension $D=D_s\oplus C$. \textbf{The complementary matrix $C$ is arbitrary; it is not required to be normal.} This relaxation is useful for lower-bounding the norm of every normal completion, but a minimizer in the relaxed class is not automatically a normal-matrix example.

\begin{lemma}[Reduction to the defect space]\label{lem:defect}
Every such completion admits a compression, followed if necessary by adding zero complementary coordinates, to the canonical form on $2p-1$ coordinates
\[
 U a_0=a_0,\qquad
 U\big|_{\operatorname{span}\{a_j,b_j\}}=
 \begin{pmatrix}\lambda_j&\sqrt{1-\lambda_j^2}\\
 \sqrt{1-\lambda_j^2}&-\lambda_j\end{pmatrix},\quad 1\le j<p.
\]
The compressed diagonal block is still $RT$, the remaining block is arbitrary, and the anticommutator norm does not increase.
\end{lemma}
\begin{proof}
In selected Fourier coordinates write $U=\left(\begin{smallmatrix}\Lambda&Y\\Y^*&Z\end{smallmatrix}\right)$. The identity $U^2=I$ gives $YY^*=I-\Lambda^2$. For each $s_j=\sqrt{1-\lambda_j^2}>0$, set $b_j=Y^*a_j/s_j$. These vectors are orthonormal. The off-diagonal identity and $U^2=I$ give
\[
 U a_j=\lambda_j a_j+s_j b_j,\qquad U b_j=s_j a_j-\lambda_j b_j.
\]
For example, the component of $Zb_j$ in the defect space is $-\lambda_jb_j$, while $\|Zb_j\|_2^2=\lambda_j^2$, leaving no orthogonal component. Coordinates with $s_j=0$ are already eigenvectors of $U$.

Thus the selected space plus the defect space reduces $U$. If $K$ denotes this subspace, then
\[
 P_K(DU+UD)|_K=D_KU_K+U_KD_K.
\]
Compression cannot increase the operator norm. The compression of $D$ is still $RT$ on the entire selected space and an arbitrary matrix on the defect space. For any missing $b_j$, add a zero coordinate to $D$ and the scalar $-\lambda_j$ to $U$. This adds a zero anticommutator block and gives the displayed form, including endpoint eigenvalues.
\end{proof}

\begin{theorem}[Exact completion norm]\label{thm:completion}
Set $\theta_j=\tfrac12\arccos\lambda_j\in[0,\pi/2]$ for $1\le j<p$. Then the minimum over all allowed completions and all complementary matrices is
\begin{equation}\label{eq:completion}
 \boxed{\inf\norm{DU+UD}=2R\max\left\{
 \cos\theta_1,\ \cos\theta_{p-1},\
 \max_{1\le j\le p-2}|\cos(\theta_j+\theta_{j+1})|
 \right\}.}
\end{equation}
The infimum is attained in dimension $2p-1$.
\end{theorem}
\begin{proof}
Use Lemma~\ref{lem:defect}. Let $W a_j=\omega^j a_j$ and $W b_j=\omega^j b_j$. This diagonal unitary commutes with the canonical $U$. Since $W^*TW=\omega^{-1}T$, the average
\[
 \Phi(D)=\frac1p\sum_{k=0}^{p-1}\omega^k W^{*k}DW^k
\]
keeps the selected block $RT$ fixed. The anticommutator norm cannot increase: every summand of its anticommutator is a unitary conjugate of $DU+UD$, multiplied by a scalar of modulus one. On the complementary coordinates, the average retains only entries $b_j\mapsto z_jb_{j+1}$, $1\le j\le p-2$. Thus it suffices for a lower bound to consider a broken forward shift.

For $j\ge1$, positive and negative eigenvectors of $U$ are
\[
 p_j=\cos\theta_j\,a_j+\sin\theta_j\,b_j,\qquad
 n_j=-\sin\theta_j\,a_j+\cos\theta_j\,b_j,
\]
with the additional positive vector $p_0=a_0$. If $P_+,P_-$ are the two spectral projections of $U$, then
\begin{equation}\label{eq:pinching}
 DU+UD=2P_+DP_+-2P_-DP_-,\qquad
 \norm{DU+UD}=2\max\{\norm{P_+DP_+},\norm{P_-DP_-}\}.
\end{equation}
The positive compression is a weighted cyclic shift; the negative compression is a weighted broken shift. Their norms are their largest edge moduli. The positive boundary edges are $R\cos\theta_1$ and $R\cos\theta_{p-1}$. On an internal edge, put
\[
 A_j=\cos\theta_j\cos\theta_{j+1},\qquad
 B_j=\sin\theta_j\sin\theta_{j+1}.
\]
The two weights are $RA_j+z_jB_j$ and $RB_j+z_jA_j$. For nonnegative $A,B$ and arbitrary complex $z$,
\begin{equation}\label{eq:edge}
 \max\{|RA+zB|,|RB+zA|\}\ge R|A-B|.
\end{equation}
When $A+B>0$, subtracting $B$ times the second expression from $A$ times the first proves this by the triangle inequality; when $A+B=0$ it is immediate. Equality is attained simultaneously on every edge by choosing $z_j=-R$. Since $A_j-B_j=\cos(\theta_j+\theta_{j+1})$, the formula follows. The same canonical completion and this broken shift attain it.
\end{proof}

The averaging step is permitted precisely because $C$ was not required to remain normal. At the displayed minimizer, $C=-RS_{p-1}$ is a nonzero nilpotent shift and is not normal. This is a norm-completion theorem, not a newly constructed normal pair.

\section{The polygon obstruction without a circulant assumption}
\begin{proposition}[Minimizing over the overlap]\label{prop:angles}
The minimum of the right side of \eqref{eq:completion}, over all allowed Fourier eigenvalues, is $\delta_p$.
\end{proposition}
\begin{proof}
Let $M$ be the maximum inside \eqref{eq:completion}, and write $M=\sin\gamma$ with $0\le\gamma\le\pi/2$. The boundary terms imply
\[
 \theta_1,\theta_{p-1}\ge\frac\pi2-\gamma.
\]
The internal terms imply
\[
 \frac\pi2-\gamma\le\theta_j+\theta_{j+1}\le\frac\pi2+\gamma.
\]
Because $p$ is odd, the alternating sum of the $p-2$ internal sums equals $\theta_1+\theta_{p-1}$. It has one more positive than negative term. Therefore
\[
 \pi-2\gamma\le\theta_1+\theta_{p-1}
 \le\frac\pi2+(p-2)\gamma,
\]
so $\gamma\ge\pi/(2p)$.

For equality, take $\gamma=\pi/(2p)$ and
\begin{equation}\label{eq:optimal-angles}
 \theta_j=\begin{cases}\pi/2-j\gamma,&j\text{ odd},\\j\gamma,&j\text{ even}.
 \end{cases}
\end{equation}
Both boundary angles are $\pi/2-\gamma$, and the internal sums alternate between $\pi/2+\gamma$ and $\pi/2-\gamma$. Every relevant modulus is $\sin\gamma$.
\end{proof}

\begin{theorem}[Arbitrary-overlap polygon obstruction]\label{thm:polygon}
Let $U$ be any finite-dimensional Hermitian unitary, with leading selected block $X$ of size $p$, where $p\ge3$ is odd. Let $D=D_s\oplus C$, with $C$ arbitrary. If $X$ has an eigenvalue $+1$ or $-1$, then
\begin{equation}\label{eq:polygon-obstruction}
 \boxed{\quad\norm{DU+UD}\ge 2R\sin\frac\pi{2p}=\delta_p.\quad}
\end{equation}
No circulant, realness, or fixed-completion-dimension assumption is required. The bound is sharp in the class with arbitrary complementary matrices.
\end{theorem}
\begin{proof}
Replace $U$ by $-U$ if necessary and choose a selected unit vector $v$ with $Xv=v$. The completion identity implies $U(v,0)=(v,0)$: its complementary component has squared norm $v^*(I-X^2)v=0$.

Let $P$ cyclically permute the selected polygon coordinates, so $P D_sP^*=\omega^{-1}D_s$, and extend it as the identity on the complementary coordinates. For $0\le g<p$, form
\[
 U_g=P^gUP^{*g},\qquad D_g=\omega^gP^gDP^{*g}=D_s\oplus\omega^g C.
\]
Each anticommutator has exactly the original norm. Their direct sum also has that norm, and fixes the vector consisting of the $p$ copies $P^gv/\sqrt p$ on the selected coordinates.

For each selected spectral value indexed by $i$, let $J_i$ have, in copy $g$, the single selected coordinate $i$ with amplitude $v_{i-g}$, with indices interpreted modulo $p$. The $J_i$ are orthonormal because
$\sum_g|v_{i-g}|^2=1$. Their span reduces the direct-sum $D$: each $J_i$ is an eigenvector with eigenvalue $(D_s)_{ii}$, and it lies in a reducing selected eigenspace. The selected compression of the direct-sum reflection in this basis is
\begin{equation}\label{eq:twirl}
 \widetilde X_{ij}=\sum_{g=0}^{p-1}\overline{v_{i-g}}\,
 X_{i-g,j-g}\,v_{j-g}.
\end{equation}
It is a Hermitian contraction. The cyclic change of summation index makes it circulant. Moreover, the fixed direct-sum vector equals $p^{-1/2}\sum_iJ_i$, so $\widetilde X\mathbf1=\mathbf1$.

In the orthogonal decomposition beginning with the $J_i$, the direct-sum pair is therefore an allowed completion of a circulant overlap, with selected block still $D_s$. Theorems~\ref{thm:completion} and Proposition~\ref{prop:angles} give the lower bound. The example in \eqref{eq:optimal-angles}, completed by the broken shift, attains it in the relaxed class.
\end{proof}

\begin{corollary}[Consequence for selected Hall witnesses]\label{cor:hall}
In a reflection anti-pair $D,-UDU$ of order $2p-1$, a selected regular odd polygon of size $p$ cannot yield a selected Hall quotient greater than one. This remains true with an arbitrary normal complementary block, or even without its normality.
\end{corollary}
\begin{proof}
At least one eigenspace of the Hermitian unitary $U$ has dimension at least $p$. It intersects the $p$-dimensional selected coordinate space nontrivially. Consequently $X$ has eigenvalue $+1$ or $-1$, and Theorem~\ref{thm:polygon} applies. Also $(D+UDU)U=DU+UD$, so the denominator is the actual perturbation norm whenever the pair is normal.
\end{proof}

This is a statement about the \emph{specified selected witness}. Other Hall witnesses of a completed pair can have different numerators; the corollary does not assert a constant-one bound for its entire matching distance.

\section{Compressed optima approaching $\pi/2$ do not transfer}
\begin{theorem}[An exact family of misleading compressed quotients]\label{thm:misleading}
For the circulant problem above,
\[
 \min_{X\mathbf1=\mathbf1,\ X=X^*,\ \norm X\le1}
 \norm{D_sX+XD_s}=\frac{2R}{p},
\]
where the minimum here is over circulant $X$. It is attained by
\begin{equation}\label{eq:compressed-optimal}
 \lambda_j=(-1)^j\left(1-\frac{2j}{p}\right),\qquad 0\le j<p.
\end{equation}
For this same overlap, the exact least full completion norm is $2R/\sqrt p$.
\end{theorem}
\begin{proof}
In the Fourier basis the compressed anticommutator is a weighted cyclic shift with weights $R(\lambda_j+\lambda_{j+1})$. Odd alternating telescoping gives
\[
 2\lambda_0=\sum_{j=0}^{p-1}(-1)^j(\lambda_j+\lambda_{j+1})=2,
\]
so its largest weight modulus is at least $2R/p$. Formula~\eqref{eq:compressed-optimal} lies in $[-1,1]$ and gives every adjacent sum equal to $(-1)^j2/p$, including the wraparound edge.

For the completed problem, both boundary cosines in \eqref{eq:completion} equal $1/\sqrt p$. On an internal edge, direct substitution into the half-angle identities gives, with $q=j(p-j-1)\ge0$,
\[
 |\cos(\theta_j+\theta_{j+1})|
 =\frac{\sqrt{q+p}-\sqrt q}{p}
 =\frac1{\sqrt{q+p}+\sqrt q}\le\frac1{\sqrt p}.
\]
The exact completion formula proves the claim.
\end{proof}

Thus the compressed selected quotient and the best completed selected quotient for the same overlap are, respectively,
\begin{equation}\label{eq:contrasting-limits}
 p\sin\frac\pi{2p}\longrightarrow\frac\pi2,
 \qquad
 \sqrt p\sin\frac\pi{2p}\longrightarrow0.
\end{equation}
The following decimals merely illustrate these exact formulas.
\begin{center}
\begin{tabular}{rrr}
\toprule
$p$&Compressed selected quotient&Best completed selected quotient\\\midrule
3&1.500000&0.866025\\
5&1.545085&0.690983\\
9&1.562834&0.520945\\
17&1.568562&0.380432\\
33&1.570203&0.273337\\\bottomrule
\end{tabular}
\end{center}
For $p=3$, \eqref{eq:compressed-optimal} recovers the preceding package's overlap with eigenvalues $1,-1/3,-1/3$; the full optimum $2R/\sqrt3$ recovers its completion obstruction. The result here proves the exact all-odd-order family and removes symmetry from the constant-one obstruction, rather than simply retesting the three-dimensional example.

\section{A stationary normal operator on two spectral lines}
The large inherited finite examples suggest a two-line bulk pattern. To investigate it without treating it as a finite certificate, work on
$\mathcal H=L^2(\mathbb R/2\pi\mathbb Z;\mathbb C^2)$ and define the normal operator
\begin{equation}\label{eq:D-infinite}
 \mathcal D=-\ii\frac{\mathrm d}{\mathrm d\vartheta}+\frac14I+\ii Y,
 \qquad Y=\diag(y_s,y_c),
\end{equation}
with domain the periodic Sobolev space $H^1$. The derivative part is self-adjoint, and it commutes with the constant self-adjoint $Y$, proving normality. The two labeled spectral ladders are
\[
 m+\frac14+\ii y_s,\qquad m+\frac14+\ii y_c,\qquad m\in\mathbb Z.
\]
Write $B=y_s+y_c\ge0$ and $C=y_s-y_c\ge0$.

Let $\phi$ be locally Lipschitz, odd, and satisfy $\phi(\vartheta+2\pi)=\phi(\vartheta)+\pi$. Define
\[
 V(\vartheta)=e^{-\ii\vartheta/2}
 \bigl(\ii\sin\phi(\vartheta)I+\cos\phi(\vartheta)\sigma_x\bigr),
 \quad \sigma_x=\begin{pmatrix}0&1\\1&0\end{pmatrix},
 \quad (Jf)(\vartheta)=f(-\vartheta).
\]
The matrix $V$ is periodic, unitary and symmetric, and $V(-\vartheta)=\overline{V(\vartheta)}$. Therefore $\mathcal U=VJ$ is a self-adjoint unitary. Since $V$ is Lipschitz, it preserves the domain needed for the following computation.

\begin{proposition}[Exact multiplier norm]\label{prop:multiplier}
The anticommutator, initially defined on $H^1$, extends to a bounded operator and satisfies
\begin{equation}\label{eq:model-norm}
 \boxed{\norm{\mathcal D\mathcal U+\mathcal U\mathcal D}
 =\esssup_\vartheta\left(\sqrt{\phi'(\vartheta)^2+B^2}
 +C|\sin\phi(\vartheta)|\right).}
\end{equation}
The same is the norm of the bounded perturbation $\mathcal D+\mathcal U\mathcal D\mathcal U$.
\end{proposition}
\begin{proof}
The derivative terms cancel against reflection of the angular variable. The anticommutator is
\[
 \left[-\ii V'+\frac12V+\ii(YV+VY)\right]J.
\]
Multiplying the matrix factor on the right by $V^*$ gives
\[
 -\ii V'V^*+\frac12I+\ii(Y+VYV^*)
 =\phi'\sigma_x+\ii\bigl(BI+C\sin\phi\,N(\phi)\bigr),
\]
where $N(\phi)=\sin\phi\,\sigma_z-\cos\phi\,\sigma_y$ is a Hermitian involution anticommuting with $\sigma_x$. For real $t,d,B$, the squared singular values of
$t\sigma_x+\ii(BI+dN)$ are
\[
 \bigl(\sqrt{t^2+B^2}\pm |d|\bigr)^2.
\]
This follows by multiplying by the adjoint, or by the quadratic characteristic polynomial recorded in the symbolic audit. The norm of a bounded multiplication operator is the essential supremum of its pointwise norms. Finally, right multiplication by $\mathcal U$ preserves the norm and gives the stated perturbation identity on the common domain.
\end{proof}

\Needspace{20\baselineskip}
\section{The phase optimization is an exact integral problem}
Fix $B,C>0$ and impose $\phi(0)=0$, $\phi(\pi)=\pi/2$. Paths on this interval need not be monotone.
\begin{theorem}[Exact optimal phase cost]\label{thm:phase}
The least essential supremum in \eqref{eq:model-norm} is the unique number $T>B+C$ satisfying
\begin{equation}\label{eq:time-integral}
 \int_0^{\pi/2}\frac{\dd u}{\sqrt{(T-C\sin u)^2-B^2}}=\pi.
\end{equation}
It is attained by the increasing solution of
\begin{equation}\label{eq:ode}
 \phi'=\sqrt{(T-C|\sin\phi|)^2-B^2},\qquad \phi(0)=0,
\end{equation}
extended to the real line. This extension has exactly the symmetries required in Section~5.
\end{theorem}
\begin{proof}
For $T>B+C$, the integral in \eqref{eq:time-integral} is continuous and strictly decreasing in $T$. It tends to zero at infinity and diverges as $T\downarrow B+C$: near $u=\pi/2$ the denominator is bounded above by a constant times the square root of $T-B-C+(\pi/2-u)^2$. Hence there is one solution.

For any path with cost strictly below a proposed $T>B+C$, the norm constraint gives
$|\phi'|\le\sqrt{(T-C|\sin\phi|)^2-B^2}$ almost everywhere. Composing $\phi$ with the increasing primitive of the reciprocal right side and integrating yields
\[
 \int_0^{\pi/2}\frac{\dd u}{\sqrt{(T-C\sin u)^2-B^2}}\le\pi.
\]
The argument uses an absolute derivative bound and remains valid for nonmonotone paths. Every admissible path has cost at least $B+C$, since it reaches $\pi/2$. Applying the displayed inequality for all larger $T$ also excludes cost exactly $B+C$ by divergence. Monotonicity of the integral proves the lower bound.

For the unique solution $T$, the positive ODE speed in \eqref{eq:ode} takes exactly time $\pi$ to reach $\pi/2$. The speed is even in $\phi$ and $\pi$-periodic. Consequently the solution is odd and advances by $\pi$ in time $2\pi$. Its pointwise cost is exactly $T$, proving attainment and the required periodicity of $V$.
\end{proof}

Normalize $b=B/T$, $c=C/T$, so $b,c\ge0$ and $b+c<1$, and define
\begin{equation}\label{eq:I}
 I(b,c)=\int_0^{\pi/2}\frac{\dd u}{\sqrt{(1-c\sin u)^2-b^2}}.
\end{equation}
Then $T=I(b,c)/\pi$. The selected-line pair-sum separation is
$\sqrt{1/4+(B+C)^2}$, giving the \emph{formal selected-gap quotient}
\begin{equation}\label{eq:formal}
 \Rformal(b,c)=\sqrt{(b+c)^2+\left(\frac{\pi}{2I(b,c)}\right)^2}.
\end{equation}
This is not yet a matching distance; Section~7 computes the actual infinite matching explicitly.

\begin{proposition}[Unique height balance at fixed sum]\label{prop:convex}
Fix $s=b+c\in(0,1)$ and vary $c\in[0,s]$. The function $I(s-c,c)$ is strictly convex. Its unique minimum is at $c=0$ when $s\le2/\pi$, and at a unique interior point when $s>2/\pi$. Thus the maximum of \eqref{eq:formal} at fixed $s$ has the same unique choice. The optimized formal quotient equals one on $s\le2/\pi$, exceeds one on $s>2/\pi$, and tends to one as $s\uparrow1$.
\end{proposition}
\begin{proof}
Put $z=\sin u$ and
\[
 Q=1-s^2+2c(s-z)-c^2(1-z^2)>0,\qquad g=s-z-c(1-z^2).
\]
The second derivative of the integrand is
\[
 \frac{\partial^2}{\partial c^2}Q^{-1/2}
 =3g^2Q^{-5/2}+(1-z^2)Q^{-3/2},
\]
whose integral is strictly positive. At $c=0$ the first derivative of the integral is
$[1-(\pi/2)s]/(1-s^2)^{3/2}$. At $c=s$ it is
$\int_0^{\pi/2}z(1-sz)^{-2}\dd u>0$. These signs and strict convexity prove the minimizer assertions. At $c=0$, $I=\pi/(2\sqrt{1-s^2})$, so the formal quotient equals one; decreasing $I$ makes it larger.

For the endpoint limit, write $v=\pi/2-u$. Factoring $Q$ gives, uniformly for $0\le c\le s$,
\[
 Q=(1-s+c(1-z))(1+s-c(1+z))\le2(1-s)+v^2.
\]
Hence $I(s-c,c)\ge\operatorname{arsinh}(\pi/(2\sqrt{2(1-s)}))\to\infty$ uniformly. Formula~\eqref{eq:formal} proves the limit.
\end{proof}
This reduces the height balance to a unique scalar minimization at each $s$, but does not prove uniqueness or a certified value of the final maximum over $s$. The reported numerical optimizer is not substituted for such a proof.

\section{A certified model number and the lost Hall obstruction}
For the fixed rational parameters
\begin{equation}\label{eq:chosen}
 b=\frac{296277}{500000}=0.592554,\qquad
 c=\frac{74553}{200000}=0.372765,
\end{equation}
the standard-library checker proves
\begin{equation}\label{eq:certified-model}
 \boxed{1.0373054176426<\Rformal(b,c)<1.0373054176428.}
\end{equation}
These inequalities concern this parameter pair only. Three local numerical searches gave values near $1.0373054176438066$, but neither that last number nor uniqueness of the parameter optimum has been certified.

The finite arithmetic proof of \eqref{eq:certified-model} is described in Appendix~A. The substitution $t=\tan(u/2)$ converts the integral to
\begin{equation}\label{eq:rational-integral}
 I(b,c)=2\int_0^1\frac{\dd t}{\sqrt{Q(t)}},\quad
 Q(t)=(1+t^2-2ct)^2-b^2(1+t^2)^2.
\end{equation}
The checker uses 32 cells, a degree-12 binomial expansion on each cell, exact integrated polynomials, and a geometric remainder bound. Square roots and accumulation are rounded outward at rational dyadic endpoints. A Machin-series enclosure supplies $\pi$. No floating-point comparison decides either bound. The width of the computed integral enclosure is less than $2\cdot10^{-23}$; the stated decimal interval deliberately leaves a larger margin.

\begin{proposition}[The actual infinite matching has ratio below one]\label{prop:actual-matching}
For the two labeled spectral ladders in Section~5 and their negatives, define the bottleneck using all bijections between the countable labeled lists. Its value is exactly
\begin{equation}\label{eq:actual-distance}
 d_{\infty,\mathbb Z}=\sqrt{\frac14+B^2}.
\end{equation}
For the optimal phase with $B,C>0$,
\[
 \frac{d_{\infty,\mathbb Z}}{\norm{\mathcal D\mathcal U+\mathcal U\mathcal D}}
 <1.
\]
At \eqref{eq:chosen}, this actual ratio is approximately $0.703762744557$, not $1.037305417643$.
\end{proposition}
\begin{proof}
Match the selected point with index $m$ to the negative complementary point with index $-m$, and conversely. Each displacement is $1/2+\ii(y_s+y_c)$, so this is a bijection of cost $\sqrt{1/4+B^2}$.

For the reverse bound, every selected point is at least this distance from every negative complementary point: their real displacement lies in $\mathbb Z+1/2$, and their imaginary displacement is $B$. Its distance to a negative selected point is at least $\sqrt{1/4+(B+C)^2}$, which is no smaller. Thus no bijection can have smaller cost.

In the phase equation, the integrand is strictly larger, except at an endpoint, than $1/\sqrt{T^2-B^2}$ because $C>0$. Therefore
\[
 \pi=\int_0^{\pi/2}\frac{\dd u}{\sqrt{(T-C\sin u)^2-B^2}}
 >\frac{\pi}{2\sqrt{T^2-B^2}},
\]
which gives $T>\sqrt{1/4+B^2}$. The numerical value is separately enclosed by the same rational integral audit.
\end{proof}

This is the exact reason that the formal selected gap cannot simply be called a matching distance. In a finite example with $p$ selected vertices and $p-1$ complementary vertices, a cross-channel matching leaves a cardinality deficit. In the two bilateral infinite ladders, the displayed cross-channel bijection has no such deficit. Passing to an infinite bulk discards this combinatorial information.

\section{Comparison with the retained finite example}
The numerical model was compared with the \emph{floating-point search candidate} retained in the prior archive, not substituted for its exact rational certificate. For the chosen parameters, let $y_s=(b+c)T/2$, $y_c=(b-c)T/2$. After normalizing the selected entry of index zero to one, the proposed spectral pattern is
\[
 a_m=1+\frac{m}{1/4+\ii y_s},\qquad
 c_m=\frac{1/4+\ii y_c+m}{1/4+\ii y_s},
\]
with indices ordered $0,-1,1,-2,2,\ldots$. This explains the nearly affine two-line bulk observed in the prior search. It is not an assertion that the finite example lies exactly on those lines.

The Fourier coefficients of the two entries of $V$ are
\[
 h_k=\frac1\pi\int_0^\pi\sin\phi(\vartheta)
             \sin((k+1/2)\vartheta)\dd\vartheta,\quad
 g_k=\frac1\pi\int_0^\pi\cos\phi(\vartheta)
             \cos((k+1/2)\vartheta)\dd\vartheta.
\]
Selected rows of the inherited overlap agree closely with these coefficients after the sign convention documented in the diagnostic source. The following are floating-point diagnostics, not certificates.
\begin{center}
\begin{tabular}{rrrrr}
\toprule
$k$&Model $h_k$&Finite $h_k$&Model $g_k$&Finite $g_k$\\\midrule
0&0.56426524&0.56423928&0.41306541&0.41309974\\
1&0.07891353&0.07891353&0.06440656&0.06442765\\
2&0.01945124&0.01945229&0.01443338&0.01444997\\\bottomrule
\end{tabular}
\end{center}
The largest spectral discrepancy among the first 21 selected coordinates is approximately $0.00148$, and among the first 21 complementary coordinates it is approximately $0.00140$. Over all coordinates, these discrepancies rise to approximately $0.246$ and $0.0449$, respectively. The boundary deformation is therefore not uniformly negligible in the recorded example.

The retained candidate's floating-point perturbation norm is about $1.92810235454$. The model norm after rescaling its selected gap to two is about $1.92807245246$. Their proximity motivates the model but proves neither convergence nor extremality. The rational $193\times193$ certificate, reverified independently of this fit, continues to establish only the unchanged lower endpoint in \eqref{eq:bracket}.

A valid finite transfer theorem would have to construct actual finite normal diagonal matrices $D_N$, Hermitian-unitary reflections $U_N$, and selected sets of size $p_N$ in dimension $2p_N-1$, with
\[
 \min_{i,j\in S_N}|(D_N)_{ii}+(D_N)_{jj}|\ge\delta-o(1),
 \qquad \norm{D_NU_N+U_ND_N}\le T+o(1).
\]
That would prove the corresponding finite lower bound $\Cstar\ge\delta/T$. A pointwise bulk fit or strong operator limit alone does not give the necessary norm upper bound at the finite boundary. No such transfer theorem is proved here. Even a successful transfer would establish a lower bound, not the missing unrestricted upper bound.

\section{Verification, scope, and the remaining theorem}
The delivered programs distinguish exact arithmetic from numerical diagnostics. The current execution record contains the following checks.

The standard-library polygon audit checks the rational eigenvalue and cyclic-sum identities, the exact angle-coefficient telescoping, and the square-root comparison identities for all 500 odd orders from 3 through 1001. These finite checks supplement the general proofs in Sections~2--4. The model checker recomputes the rational integral and the strict comparisons in \eqref{eq:certified-model}. Ten exact regression and rejection-test methods pass, including 1,200 square-root enclosure cases, deliberately false claimed intervals, invalid parameters, and repeat execution with Python optimization enabled.

Eight numerical regression groups pass. They include 120 canonical attainment cases, 150 arbitrary complementary matrices, 250 random angle sets in addition to 25 exact-formula specializations, 25 compressed-optimum examples, 120 noncirculant orbit compressions, 15 explicit orbit dilations, and 600 multiplier-norm checks. These are finite diagnostics and do not certify a universal inequality by themselves. A separate SymPy reconstruction passes 13 checks, including the full $5\times5$ radical examples, their exact Gram characteristic polynomials, the Pauli singular-value identity, and the strict-convexity derivative identity.

The retained rank-one arithmetic certificate and its nine exact rejection tests pass again. The inherited $193\times193$ normal pair passes both original exact positivity routes: congruence diagonal dominance and outward-rounded interval $LDL$ elimination. The two routes share parts of the enclosure code and are not an independent external mathematical review.

From the extracted archive root, the new exact checks require only Python's standard library:
\begin{verbatim}
python -S verify_polygon.py
python -S verify_model.py
python verify_all.py
\end{verbatim}
Optional numerical, symbolic, and inherited checks are run with
\begin{verbatim}
python verify_all.py --numerical --prior
\end{verbatim}
The prior ZIP is retained unchanged. Source code, chosen-parameter data, the inherited search candidate used for the fit, executed logs, optional dependency versions, and a file-hash manifest are included. None of the numerical optimizer stopping flags is treated as a proof of a global maximum.

\textbf{The unrestricted gap.} Using the inherited real-reflection reduction, a full solution still needs an explicit $L$ and the inequality
\begin{equation}\label{eq:missing}
 \min_{1\le i,j\le k+1}|a_i+a_j|
 \le L\norm{DU+UD}
\end{equation}
for every $k$, every complex diagonal $D=\diag(a_1,\ldots,a_{2k+1})$, and every real symmetric orthogonal $U$ of inertia $(k+1,k)$, together with constructions attaining or approaching the same $L$. Section~3 proves a constant-one statement for a particular geometry, not for arbitrary diagonal data. Sections~5--7 optimize a stationary phase model, not the supremum in \eqref{eq:missing}.

The new work rules out a broad compressed-polygon shortcut and identifies exactly why the apparent stationary quotient is not an infinite matching ratio. It does not determine $\Cstar$, the exact rank-two constant, a uniform rank bound, or the extremal geometry. The repository status should not be changed to solved on the basis of this package \cite{guidelines}. No repository content was modified. Historical priority, external peer review, and proof-assistant verification are not claimed.

\appendix
\section{Exact quadrature certificate}
For the rational $b,c$ in \eqref{eq:chosen}, expand \eqref{eq:rational-integral} as
\[
 Q(t)=(1-b^2)-4ct+(2+4c^2-2b^2)t^2-4ct^3+(1-b^2)t^4.
\]
On a cell centered at $m$ with half-width $h$, put $t=m+hy$ and write
\[
 Q(m+hy)=q_0(1+u(y)),\qquad
 u(y)=u_1y+u_2y^2+u_3y^3+u_4y^4.
\]
Every coefficient is rational. The checker verifies $q_0>0$ and
$r=\sum_{j=1}^4|u_j|<1$. For $|y|\le1$, the binomial expansion has coefficients
$c_k=(-1)^k\binom{2k}{k}/4^k$, with $|c_k|\le1$. Consequently
\[
 \left|(1+u)^{-1/2}-\sum_{k=0}^{K}c_ku^k\right|
 \le\frac{r^{K+1}}{1-r}.
\]
Integrating the polynomial exactly over $[-1,1]$ uses only its even coefficients, since $\int_{-1}^1y^{2j}\dd y=2/(2j+1)$. The integrated error is at most $2r^{K+1}/(1-r)$. Multiplication by $2h/\sqrt{q_0}$ gives a rigorous cell enclosure for $I$.

For a positive rational $q=n/d$ and scale $S=2^b$, the integer
$k=\lfloor\sqrt{\lfloor nS^2/d\rfloor}\rfloor$ satisfies
$k/S\le\sqrt q\le(k+1)/S$; equality is checked separately. This supplies outward rational square-root and reciprocal bounds. Cell endpoints are rounded outward before summation, avoiding an unbounded product of unrelated rational denominators. The delivered calculation uses $K=12$, 32 cells, and 128-bit output endpoints, with additional internal square-root bits.

For $\pi$, use the alternating series for $\arctan x$ with its next-term remainder and Machin's identity
\[
 \pi=16\arctan(1/5)-4\arctan(1/239).
\]
The tangent addition identity is checked exactly: $\tan(4\arctan(1/5))=120/119$, and subtracting $\arctan(1/239)$ gives tangent one in the first quadrant. The alternating-series intervals therefore give rational lower and upper bounds for $\pi$.

The recomputed enclosure also yields the simpler rational interval
\[
 \frac{413708879761528711367}{10^{20}}<I(b,c)<\frac{413708879761528711404}{10^{20}}.
\]
Together with $314159265358979323846/10^{20}<\pi<314159265358979323847/10^{20}$, these coarser rational endpoints alone suffice for the advertised model-ratio comparisons; the checker verifies both squared comparisons by exact cross-multiplication.

Finally, monotonicity in the positive quantities gives
\[
 (b+c)^2+\left(\frac{\pi_-}{2I_+}\right)^2
 \le\Rformal(b,c)^2\le
 (b+c)^2+\left(\frac{\pi_+}{2I_-}\right)^2.
\]
A last outward square-root enclosure proves the stated strict decimal comparisons. Replacing $(b+c)^2$ by $b^2$ encloses the actual infinite matching ratio from Proposition~\ref{prop:actual-matching}. The verifier recomputes these inequalities rather than trusting stored interval endpoints or a success flag.

\begin{thebibliography}{9}
\bibitem{sp07} Open Problems in Numerical Linear Algebra.
\emph{SP-07: The sharp spectral-matching constant for normal matrices.}
Canonical public problem statement, read during this continuation.
\url{https://raw.githubusercontent.com/ajt60gaibb/OpenProblemsInNLA/main/eigenvalues-and-inverse-problems/SP-07/README.md}.

\bibitem{tang} Qiyue Tang.
\emph{A Truncated Singular-Value Bound for Spectral Variation of Normal Matrices.}
arXiv:2609.09177v1, introduction and Section 3. The HTML and the PDF introduction were checked; the quoted universal upper bound is reported there.
\url{https://arxiv.org/html/2609.09177v1}.

\bibitem{prior4} \emph{SP-07: A sharp rank-one-reflection theorem.}
Preceding continuation supplied in this conversation. Its unchanged archive is retained as
\path{prior/SP-07_round4_sharp_subclass.zip}. The earlier real-reflection reductions and the exact $193\times193$ certificate remain nested inside that archive. New results are not attributed to an external audit of these prior packages.

\bibitem{guidelines} Open Problems in Numerical Linear Algebra.
\emph{Contribution and status guidelines.} Read during this continuation.
\url{https://raw.githubusercontent.com/ajt60gaibb/OpenProblemsInNLA/main/CONTRIBUTING.md}.
\end{thebibliography}
\end{document}
